\appendix
\FloatBarrier
\section{Reading glossary}\label{sec:glossaire}
Entries follow the order of first appearance in the body of the
text, after the abstract. Symbols introduced together are grouped;
their operative definitions remain those of the corresponding sections.

\begin{small}
\begin{longtable}{@{}>{\raggedright\arraybackslash}p{.27\linewidth}p{.66\linewidth}@{}}
\toprule Term or notation & Meaning in this article \\
\midrule\endhead
Rebound effect & Deviation from the consumption savings expected after an
efficiency gain. Here, the experiment measures a non-proportional
productive response, without calibrating a final consumption. \\
Endogenous; emergence & Produced by the interactions and evolution of
the system, rather than imposed as a macroscopic target. \\
Self-organized criticality (SOC) & Hypothesis of a regime spontaneously
maintaining itself near a critical threshold. It is not demonstrated
here. \\
Entity & Generic producing agent, able to lend, borrow and disappear;
none of these roles is predefined. \\
$K_i$; productive stock & Mobilizable physical capital of entity $i$,
in model joules; distinct from its net worth. \\
Technology $f_i(K)=A_iK^{\gamma_i}$ & Production added to the stock
during a step; $A_i$ is the extraction coefficient, $\gamma_i$ the
production exponent. \\
$K^\circ$, $\gamma^\circ$ & Birth endowment and production exponent
at birth in the homogeneous reference regime. \\
Principal $q$; service $rq$ & Nominal amount of a loan and interest
due per step; the principal is not amortized. \\
$C_i$, $D_i$, $NW_i$ & Claims, debts and net worth:
$NW_i=K_i+C_i-D_i$. \\
$\lambda$, $\delta$, $\sigma$, $\rho$ & Birth intensity, fraction
depreciated per step, amplitude of the log-normal shock and intensity
of market encounters. \\
Credit institution & Rules fixing the direction, amount and rate of
loans; it does not here designate a particular bank. \\
Invariant; real balance sheet & Relation imposed by the model's
operations. The real balance sheet tracks contributions, production
and stock losses. \\
$B_t$, $\eta_i$ & Number of births at step $t$ and individual
logarithmic shock, renewed at each step. \\
Concavity & Decrease of the marginal return with the stock, for
$0<\gamma<1$; origin of the possible reallocation gain. \\
Service default & Inability to fully pay the interest due at the
step considered. \\
Stock; flow; $\Delta t$ & Quantity held; quantity transferred or
produced during a step; possible physical duration of that step. \\
Exergy & Capacity of a resource to supply work relative to a
reference environment; its correspondence with economic entitlements
remains a hypothesis here. \\
$K_{\rm aut}$, $x_i$ & Deterministic fixed stock of an isolated
entity with no noise or credit; reduced capital $x_i=K_i/K_{\rm aut}$. \\
$\kappa^\circ$, $a$ & Reduced endowment $K^\circ/K_{\rm aut}$;
alternative coefficient $a=A(K^\circ)^{\gamma-1}$. \\
Rescaling; compensation & The former transforms the whole state and
its units; the latter here adapts the endowment of future entrants
to preserve $\kappa^\circ$ after the rise in $A$. \\
Tension $T$; $K_{\rm eq}$ & Ratio $K_{\rm aut}/K_{\rm eq}$; fictitious
stock giving the average production per producer. In the protocols,
$T$ can also denote the horizon of a simulation. \\
$Y$, $N_{\rm prod}$ & Total production of the step and number of
entities present in the productive phase. \\
Marginal return $f'(K)$ & Additional production at the margin per
additional unit of stock. \\
$L$, $\Delta$, $p$ & Productive loss of the lender, joint productive
surplus and share of that surplus attributed to the lender:
$rq=L+p\Delta$. \\
Bankruptcy; cascade; avalanche & Removal of an entity in default or
insolvent; propagation of losses; weakly connected component of
losses among entities that died in the same step. \\
Root; generation; $b_1$, $b_2$ & Initial death; rank of intra-step
propagation; fraction of non-root deaths and average number of
causal links between successive generations per death. \\
Seed; pairing; 95\% CI & Initialization of the random generator;
branches originating from the same state; 95\,\% confidence interval
computed across repetitions, unless stated otherwise. \\
$k$; stationarity & Historical size of the encounter pool, fixed at
two in the main protocol; temporal invariance of statistical
properties, not demonstrated by a single drift test. \\
Little's law & Relation between mean population, arrival flow and
mean lifetime under conditions of stability and boundary control. \\
$S$, $\alpha$, $S_c$, $Z$ & Avalanche size, exponent, exponential
cutoff and normalization constant of its fitted law. \\
Bootstrap; $R^2$ & Resampling of whole realizations; descriptive
goodness-of-fit score, not proof of a family of laws. \\
Contraction; $D$, $\beta$ & Sequence of decreasing steps of a
quantity; duration of that sequence and decay rate of its
exponential survival. $D$ without subscript in this context is not
a debt. \\
Gini; Lorenz & Synthetic inequality index; curve of the cumulative
share of a quantity as a function of the share of entities ranked
by it. \\
Tail; Pareto; log-normal & Domain of large values; power-law
decaying family; family whose logarithm is normal. \\
Finite elasticity $\varepsilon$ & Ratio of the logarithm of the
productive response to the logarithm of the technical factor $m$,
here $m=1.5$. \\
Quantile; decile; cohort & Ranking threshold; group of 10\,\%
reranked at each date; set of individuals tracked by identity. \\
Covariance & Measure of joint variation used in the sharing weighting
identity; it does not demonstrate causality. \\
Load ratio & Interest paid or due, depending on the definition,
relative to production plus interest received; distinct from a
debt-stock-to-production ratio. \\
Vuong & Comparison of model likelihoods on the same sample; its
assumptions and limitations are made explicit in the appendix. \\
\bottomrule
\end{longtable}
\end{small}

\clearpage
\section{Checking the regime studied and the response calculations}\label{sec:stationnarite}
For each seed $s$, arm and observable $X$, we compute
\[
 R_s(X)=\frac{\langle X\rangle_{3750<t\le4000}}
                  {\langle X\rangle_{3000<t\le3250}},\qquad
 t_X=\frac{\overline R(X)-1}{s_R/\sqrt{12}},
\]
where $s_R$ is the standard deviation of the twelve ratios. The 95\% CI is
$\overline R\pm t_{0.975;11}s_R/\sqrt{12}$, with
$t_{0.975;11}\simeq2.201$.
This diagnostic compares two distant temporal levels; it analyzes
neither the full correlation structure nor the stability of the distribution.

\begin{table}[htbp]\centering\small
\caption{Ratio of the means of the last and first 250 steps of the terminal window.}\label{tab:stationnarite}
\begin{tabular}{@{}lrrr@{}}
\toprule
Arm & $R(Y)$ (95\% CI) & $R(N)$ (95\% CI) & $R(\sum K)$ (95\% CI) \\
\midrule
Control & $1.0029\pm0.0152$ & $1.0014\pm0.0132$ & $1.0043\pm0.0178$ \\
$A\times1.5$ & $0.9968\pm0.0139$ & $0.9979\pm0.0116$ & $0.9957\pm0.0167$ \\
$A\times1.5$, compensated endowment & $0.9987\pm0.0132$ & $0.9974\pm0.0122$ & $0.9998\pm0.0148$ \\
Entrants $A\times1.5$ & $1.0088\pm0.0106$ & $1.0070\pm0.0090$ & $1.0092\pm0.0128$ \\
Entrants $A\times0.75$ & $0.9906\pm0.0159$ & $0.9928\pm0.0139$ & $0.9905\pm0.0184$ \\
Entrants $\gamma=0.6$ & $1.0097\pm0.0159$ & $1.0038\pm0.0133$ & $1.0100\pm0.0186$ \\
\bottomrule\end{tabular}\end{table}


None of these contrasts exceeds the usual two-sided threshold, but the
intervals allow drifts on the order of a percent. The temporal observations
of a single run are not used as 1000 independent repetitions.
The result does not automatically hold for an earlier window,
another observable or an unpresented exploratory configuration.

\paragraph{Renewal and memory of the initial state.}
The renewal of a decile provides another time scale, but
the ranking quantity must be specified and the disappearance of an
entity distinguished from its exit from the decile. For a set $D(t_0)$
defined at a reference date, retention
$|D(t_0)\cap D(t_0+h)|/|D(t_0)|$ measures persistence in the group,
while the fraction of its members still alive measures their survival.
These two curves are not interchangeable; the term "first
decile" must specify whether it denotes low or high values.
Even low retention does not demonstrate forgetting of the
initial configuration: commitments and correlations may persist beyond
individuals. This check must be confronted with the autocorrelations of
the observables and with distinct initializations. No quantified
forgetting time is established here; the distance from the final window and
the absence of detected drift are not enough to certify it.

Population balances give exactly
$N(t+1)-N(t)=B(t)-D(t)$ if $B$ and $D$ denote the entries and exits
of the step under the same counting convention. Over a long stable window,
$\overline B\simeq\overline D$ and $\overline B\simeq\lambda$;
the ratio $\overline D/\overline N$ becomes close to
$\lambda/\overline N$. This is not exactly the mean of
$D(t)/N(t)$, and boundary terms are not zero by decree.
The relation explains why a mortality elasticity can be
almost the opposite of a population elasticity without constituting an
independent causal finding.

The elasticity of equation~\eqref{eq:eps} concerns a finite
intervention. Its mean and CI are computed after forming the twelve
paired log-ratios, not after averaging all time points
in both arms. The decomposition formulas are sensitive to the
choice of population: entities present at production may
disappear before the end-of-step snapshot. This difference is
preserved in the source data; it must not be erased to
force the numerical closure of an ill-defined identity.

\section{Distribution fits: definitions and limitations}\label{sec:ajustements}
\subsection{Mass, per-bin density and survival function}
For an integer size $S$, the mass $P(S=s)$ and the survival
$P(S\ge s)$ are distinct. The histogram here represents the mass
of a bin $[l,r[$ divided by $r-l$, so that a wide bin does not
appear more probable merely because it contains more
sizes. The fitted curve is summed over the same integer sizes
before this division. A survival slope cannot be
substituted directly for the mass exponent.

The descriptive likelihood of the law~\eqref{eq:avalanche} is maximized
per seed, over all sizes $s\ge1$. The numerical normalization
uses $1\le s\le20000$; the mass of the last thousand values is
checked to be below $10^{-10}$ at the optimum. The numerical search
bounds are $-2\le\alpha\le5$ and
$10^{-5}\le1/S_c\le2$; an optimum at a bound would signal
insufficient identification. The 24 fits presented are interior.
The fitting support is not chosen to maximize the apparent
straightness of a subset of points.

The cross-seed CIs and the bootstrap over whole runs measure the
variability of the estimators or of the average curves obtained by this
protocol. Fitting by product of probabilities does not demonstrate
the independence of a run's avalanches. We therefore do not derive
event-level standard errors or a certified goodness-of-fit test from it.
In particular, this should be complemented by a comparison of families on
the same support and an absolute-adequacy diagnostic accounting
for temporal dependence \cite{clauset}.

The descriptive coefficient shown is
\[
 R^2_{\log}=1-
 \frac{\sum_j[\log\widehat p_j-\log p_{j,\rm fit}]^2}
      {\sum_j[\log\widehat p_j-\overline{\log\widehat p}]^2},
\]
over the non-empty empirical bins. This is not a percentage
of correctly predicted probability, and it is not the criterion
optimized by the likelihood. It depends on the binning. Zeros are
not replaced by an arbitrary small value to make them
loggable.

\subsection{Durations: discrete exponential and window choice}
If $D\ge1$ follows a geometric law, with stopping probability $u$ at
each step, then $P(D\ge d)=(1-u)^{d-1}$ and
$\widehat u=1/\overline D$. Thus
$\widehat\beta=-\log(1-1/\overline D)$ in
$P(D\ge d)=\exp[-\beta(d-1)]$.
The regressions of figure~\ref{fig:art_cycles} use this
normalized form, not a log-log line. The $R^2$ there is computed
on the non-zero log-survivals, at durations 1 to 11. Episodes remain
potentially correlated and the window boundaries truncate some
of them; the CIs are therefore based on the five seeds and not on
a fictitious number of independent recessions.

\subsection{The Vuong test and what it actually compares}
In the existing analyses of net worth and interest, the threshold
$x_{\min}$ is chosen by minimizing a Kolmogorov--Smirnov
distance for a continuous power law. Above this same
threshold and on the same values, a log-normal and an
exponential are also fitted, conditioned on $X\ge x_{\min}$.
The power law compared here has \emph{no upper cutoff}: this
is not the discrete model with exponential cutoff used for avalanches.
The \emph{left} truncation, which defines a tail sample,
must not be confused with the \emph{right} cutoff.

For the fitted densities $f$ and $g$, define
$d_i=\log f(x_i)-\log g(x_i)$ and then
\[
 z=\frac{\sqrt n\,\overline d}{s_d},\qquad
 s_d^2=\frac1{n-1}\sum_i(d_i-\overline d)^2.
\]
A positive value favors $f$ in average log-likelihood.
The usual interpretation of $\pm1.96$ as a 5\%
threshold requires regularity, model-separation
and sampling assumptions. Here the threshold was selected, entities
interact, and some fitted log-normals are very close
to a power-law limit. The score is therefore reported as a diagnostic,
not as certification of these assumptions \cite{vuong}.

The flag $\sigma_{\rm LN}>10$ concerns 50 of the 72 interest
fits and 46 of the 72 net-worth fits. It marks a
numerical zone, not a mathematical boundary between two families.
Two visible groups in a scatter of fit parameters do not
mean that the agents form two economic populations:
each point in that scatter represented one fit, not one entity.
Per-fit results and their identifiers are available in
\chemin{data_revision/rev_vuong_points.csv}.

Comparing Pareto, cutoff power law and log-normal requires
fixing the same domain, describing the threshold selection and evaluating
goodness of fit, not merely counting parameters. A cutoff
that diverged with system size would define an interesting
limit, to be checked across several sizes; it does not make
its finite-size fit free of charge.

\section{Credit load and surplus weighting}
\subsection{Definition and measurement of the load ratio}\label{sec:charge}
Let $P_i$ denote the step's production, $I_i^{\text{received}}$ the interest
received, $I_i^{\text{paid}}$ that actually paid, and $D_i$ the
sum of principals due. Three ratios must be distinguished:
\[
 c_i^{\text{paid}}=\frac{I_i^{\text{paid}}}{P_i+I_i^{\text{received}}},\qquad
 c_i^{\text{due}}=\frac{\sum_{\ell:i\,\mathrm{borrows}}r_\ell q_\ell}
                         {P_i+I_i^{\text{received}}},\qquad
 d_i^{\text{income}}=\frac{D_i}{P_i+I_i^{\text{received}}}.
\]
The first is a share of gross income paid; the second is the
contractual charge due; the third relates a debt stock
to a per-step income and is interpreted in a number of steps. The
variable $D_i/P_i$ represents neither of the first two ratios.
The denominator includes financial income, even though its
aggregate sum is not new real production.

\figurine{art_charge}{0.95}{\textbf{A correctly named load
diagnostic.} Marginal rule M4.4, twelve seeds. For the snapshots
$3000<t\le3990$, the variable is
$I_i^{\text{paid}}/(P_i+I_i^{\text{received}})$; the response is the fraction
of entities dying in the following ten steps. The bins have the
explicit bounds shown, common to all seeds; they are
not capital classes. The last incomplete window is
excluded. The 95\% CIs compare the fractions obtained across twelve runs,
without attributing binomial independence to repeated entities.
All-cause deaths and insolvency roots are shown separately.}

The panels record the interest paid, not the service due
at the exact moment preceding payment. A pro-rata payment can
make an unsustainable load look low; the moment of
collection and the flows received within the step also matter.
The diagram remains descriptive, conditional on the survivors observed
in the panel. It justifies neither a causal link from load to
death nor a general refutation of convexity. Convexity
calculations or Jensen-gap computations on $D_i/P_i$ would not answer
the question asked here about the load ratio as defined.

\subsection{Proof of the weighting identity}\label{sec:ponderation}
For $n$ contracts whose mean surplus is strictly positive,
let $\overline p=n^{-1}\sum p_j$,
$\overline\Delta=n^{-1}\sum\Delta_j$ and
$p_w=\sum p_j\Delta_j/\sum\Delta_j$.
Defining the empirical covariance with denominator $n$,
\[
 \operatorname{Cov}_n(p,\Delta)
 =\frac1n\sum_j(p_j-\overline p)(\Delta_j-\overline\Delta)
 =\frac1n\sum_jp_j\Delta_j-\overline p\,\overline\Delta,
\]
we immediately obtain
\begin{equation}
 p_w-\overline p=\frac{\operatorname{Cov}_n(p,\Delta)}{\overline\Delta}.
 \label{eq:identite}
\end{equation}
If a corrected covariance with denominator $n-1$ is used,
the numerator must be multiplied by $(n-1)/n$. The identity also
holds in expectation when the moments exist and
$\mathbb E[\Delta]\ne0$.

For a constant fixed sharing rule, the gap is zero by construction.
For a variable rule, it can be positive even if most
contracts have a very small principal or surplus. The
effective distribution of contracts determines its magnitude; an
analytical grid covering rarely realized contracts cannot
replace this distribution. Qualifying a rule from
these theoretical cases alone would overstate their empirical reach.
The accounting identity does not demonstrate an explanation of
mortality or population differences between institutions.

\section{Measurement supplements and tabulated values}\label{sec:tableaux}
\subsection{Quantile ratios}
The table below and figure~\ref{fig:art_quantiles} rest
on the same observations. Each $\pm$ is a half 95\% Student CI
of the twelve per-seed ratios, not a standard deviation across
entities. "Gross income" means production plus interest received.
\begin{small}\input{table_quantiles_en}\end{small}

\subsection{Temporary reach of an intervention}
An intervention limited to the entities present can die out
with their cohort if entrants do not inherit it. This
partial-reach experiment is a protocol control,
not a main proof of the rebound mechanism. Conversely,
modifying all entities and future entrants preserves
the treatment despite population renewal. Modifying
$\gamma$ changes the shape of the function and the dimensions of
its coefficient: this is not the same operation as multiplying
$A$ at fixed $\gamma$.

\subsection{Multiple losses and sufficient causes}
The graph can attribute several losses to a victim.
This does not mean that each would have sufficed, in isolation, to
make it insolvent. To illustrate this, an initial net worth
of 5 and two losses of 3 lead to $-1$; no single loss
suffices. With a net worth of 2, each of the two losses
would suffice. It is the solvency margin before the losses that
makes the difference, not merely the number of edges.

The count of parents must also specify whether it retains only
losses coming from the previous generation or all
those of the step. A detailed sufficiency analysis also requires
reconstructing the debts written off; it is not
necessary to establish the existence of propagation nor the
non-proportionality of the response. Earlier calculations
remain archived, but their percentages are not used
as evidence in this article.

\subsection{Stock and production on the same realization}
\figurine{rev_stock_trajectoire}{0.95}{\textbf{Observing the difference between
stock and flow.} M4B central, seed 11, excerpt $3000<t\le3200$.
Each line follows a single realization, with no measurement
uncertainty artificially added; the orange points mark
the down-steps. The contraction durations are computed separately
on the two observables in figure~\ref{fig:art_cycles}.}

\FloatBarrier
\section{Local sources, configurations and reproducibility}\label{sec:sources_revision}
The PDF can be compiled with only the LaTeX files and figures
in this folder, without the internship report or the questionnaires.
The new calculations are produced by
\chemin{../scripts/revision_article.py}, which writes only into
this article's folder and launches no simulation.
The numerical exports, per-seed observations and
SHA-256 checksums of the sources are in \chemin{data_article/}.
The retained supplements from the previous version are in
\chemin{data_revision/}. Recomputing from the microdata
requires the local campaigns, which are not all version-controlled.

\subsection{Three protocols, samples not merged}
\begin{small}
\begin{longtable}{@{}p{.18\linewidth}p{.75\linewidth}@{}}
\toprule Family & Configuration and data \\\midrule\endhead
M4B central & $A=1$, $\gamma=1/2$, $\lambda=30$, $\delta=0.05$,
$\sigma=0.25$, $K^\circ=25$, $k=3$, geometric target. Five seeds
11--15; contractions on $1000<t\le4000$, Gini on the final state.
Simulation Lab identifiers are in the episode and confirmation
exports. The Gini sweep shown varies $\sigma$ only. \\
Live-v1 recalibration & Reference $A=1$, $\gamma=1/2$, $\lambda=30$,
$\delta=\sigma=0.01$, $K^\circ=25$, $\rho=1$; seeds 0--2,
empty start. Window $1200<t_{\rm old}\le2000$.
The fine-depreciation arm uses its own reference $\delta=0.002$.
The folders are named in \texttt{temps\_graines.csv}. \\
M4.4 main & Same physical reference as Live-v1; arithmetic target
in the homogeneous regime, free loan direction, marginal rule. Twelve seeds
0--11, intervention at step 2001, window $3000<t\le4000$.
Compensated endowment: $K^\circ=56.25$ for new entrants
after the intervention; existing stocks unchanged. \\
M4.4 sharing & $p\in\{0;0.25;0.5;0.75;1\}$ or marginal rule;
twelve seeds, same window. Ratios to the $p=0$ arm are paired.
Load diagnostic: marginal rule only, $3000<t\le3990$,
ten-step death horizon. \\
M4.4 encounters & $\rho\in\{0.5;1;1.5;2;3\}$; twelve seeds
per value. Balance-sheet Ginis are recomputed on the panels of
$3000<t\le4000$, like the windows of the other measurements. \\
\bottomrule
\end{longtable}
\end{small}

\subsection{Each figure and its data-point files}
The CSV files with no prefix and \chemin{diagnostic_chaine.json} are
relative to \chemin{data_article/}. The TeX files and paths
starting with \chemin{data_revision/} are relative to the LaTeX folder.
The manifest links the exports to the original
sources and their checksums; it does not turn an aggregated export
into complete microdata.
\begin{footnotesize}
\begin{longtable}{@{}p{.09\linewidth}p{.48\linewidth}p{.36\linewidth}@{}}
\toprule Fig. & Data / local file & Nature of the calculation \\
\midrule\endhead
\ref{fig:exemple} & \chemin{exemple_article.tex} & Four fictitious entities,
real loan and nominal recording; no simulated measurement. \\
\ref{fig:art_technologie} & \chemin{technologie_points.csv},
\chemin{technologie_courbes_v2.csv} & Analytical technologies
$2K^{1/4}$ and $K^{1/2}$; no regression. \\
\ref{fig:cascade} & \chemin{cascade_article.tex} & Fictitious five-entity
balance sheet; four-bankruptcy tree checked separately. \\
\ref{fig:art_avalanches} & \chemin{avalanches_counts.csv},
\chemin{avalanches_fits.csv}, \chemin{avalanches_points.csv} & Integer
sizes in the avalanche CSVs; fits and cross-seed CIs. \\
\ref{fig:art_cycles} & \chemin{cycles_fits.csv}, \chemin{cycles_points.csv};
episodes in \chemin{data_revision/rev_cycles_episodes.csv} & Five
M4B runs; exponential coefficient from the mean duration. \\
\ref{fig:art_lorenz} & \chemin{lorenz_points.csv}, \chemin{lorenz_ginis.csv} &
Control panels; exact dates 2010, 3000, 4000. \\
\ref{fig:art_sensibilite} & \chemin{m4b_sensibilite_points.csv},
\chemin{m4b_nw_confirmation.csv} & NW Gini on five independent final
snapshots per cell. \\
\ref{fig:art_rho} & \chemin{rho_points.csv}, \chemin{rho_fits.csv},
\chemin{rho_nw_graines.csv} & Three estimators from the campaign table,
plus recomputed NW Gini; per-seed regressions. \\
\ref{fig:art_temps} & \chemin{temps_graines.csv}, \chemin{temps_points.csv} &
Live-v1 series rescaled to a common duration and a common flow per
old step. \\
\ref{fig:art_reponse} & \chemin{elasticites_graines.csv}; curves in
\chemin{data_revision/rev_reponse_courbes.csv} & Elasticities recomputed
from the series; centered 101-step smoothing kept for the curves. \\
\ref{fig:chaine} & \chemin{corps_article.tex}; diagnostics
\chemin{diagnostic_chaine.json} & Interpretive diagram,
not a causal graph identified by intervention on each mediator. \\
\ref{fig:art_quantiles} & \chemin{data_revision/rev_quantiles_graines.csv} &
Per-seed ratios kept; abscissa transformed to logit. \\
\ref{fig:art_deciles} & \chemin{deciles_v2_points.csv},
\chemin{deciles_v2_graines.csv} & Shares of production, interest and NW;
ranked by capital at each snapshot. \\
\ref{fig:art_institutions} & \chemin{institutions_points.csv},
\chemin{institutions_fits.csv}, \chemin{institutions_references.csv} &
Ratios to $p=0$ and descriptive linear regressions. \\
\ref{fig:art_charge} & \chemin{charge_graines.csv}, \chemin{charge_points.csv},
\chemin{charge_resume.csv} & Interest paid / gross income; deaths
observable throughout the horizon. \\
\ref{fig:rev_stock_trajectoire} &
\chemin{data_revision/rev_stock_trajectoire_points.csv} &
Raw M4B excerpt; figure kept. \\
\bottomrule
\end{longtable}
\end{footnotesize}

The new age-composition checks are in
\chemin{ages_graines.csv}. The ranking conventions are
explicit in the script; ties are broken by stable order. The revision
plan and the decision register are outside the LaTeX folder, in
\chemin{../revision_2026-09-15/}, to keep the corrections to be
reported in the internship report separate.

The supplement on instantaneous median ages is computed by
\chemin{../scripts/ages_v2.py}. Its per-seed values, its averaging
convention and the checksums of the source panels are kept in
\chemin{../revision_2026-09-17/ages_v2_graines.csv} and
\chemin{../revision_2026-09-17/ages_v2.json}. It runs no
simulation and modifies no figure generator.


The detailed configurations and the M4.4 identifiers of the
initial campaign remain in
\chemin{data_revision/runs_m44.csv}, which distinguishes the
recorded parameters from the interventions. Some seeding metadata
does not record the parameters; these gaps are not filled
by silent deduction.
The technical \texttt{K0} codes from the source files are
kept for their traceability; in the scientific notation,
the endowment is now read $K^\circ$ and its reduced ratio
$\kappa^\circ$. Renaming the notation does not modify the engines.

\subsection{Engine compatibility and rule changes}
\label{sec:compatibilite}
Three levels of control must be distinguished: preserving the
functionalities of the simulation lab, reproducing a trajectory
in a reference regime, and comparing two models whose rules
differ. Only the second bears directly on dynamic parity.

The M4.2 report, section "M4B Parity", documents tests at
$A=1$, $\gamma=1/2$ and $k=2$: identity of discrete variables
and comparison of floats at relative tolerance $10^{-9}$.
The reported 1000-step probe shows no discrete divergence,
with a maximum relative deviation of the series of $3.1\,10^{-12}$.
This finding is not a bit-for-bit identity nor a guarantee at every horizon.

For M4.4, in the campaign folder
\chemin{m4_4_rebond_credit_soc/}, the file
\chemin{results/analysis/parity_deviations_8000.csv}
holds 8000 maximum deviations, all zero, over the 26 columns compared to
the reference M4.3 run, in the homogeneous regime with arithmetic
target. The test \chemin{tests/test_parity_m4_3.py} defines this
scope. The test \chemin{tests/test_v1_equivalence.py} separately compares
Live-v1 to the compatibility mode
\texttt{richest\_lends}, with a heterogeneous intervention; a negative
control checks that freeing the loan direction can change the trajectory.
The presence of this test defines a verifiable requirement; it does not by
itself replace an execution report for every version.

The M4.4 campaign paths are kept as provenance identifiers. Switching
to a different market rule must be distinguished from an instrumentation
extension, without renaming the archives the data come from. Non-regression
therefore does not allow merging the M4B, Live and M4.4 samples.

\subsection{Why temporal aggregation is not a symmetry}
\label{sec:composition_temps}
A counterexample appears even before introducing the market. With no noise,
credit or births, let $F(K)=(1-\delta)(K+AK^\gamma)$.
Two successive steps give
\[
F(F(K))=(1-\delta)^2(K+AK^\gamma)
 +(1-\delta)A\big[(1-\delta)(K+AK^\gamma)\big]^\gamma.
\]
An aggregated step, with $A'=2A$ and
$1-\delta'=(1-\delta)^2$, would instead give
\[
G(K)=(1-\delta)^2(K+2AK^\gamma).
\]
In general $F\circ F\ne G$: the second production is computed on
an already modified stock. For example, at $K=100$, $A=1$,
$\gamma=1/2$ and $\delta=0.01$, the values are approximately
$118.142$ and $117.612$. The market adds an alternation of
encounters, contracts, interest and bankruptcies; composing the Poisson
laws or the shocks alone does not compose these coupled operations.

The Live-v1 trials of figure~\ref{fig:art_temps} measure the collective
reach of this difference. After full conversion of the parameters,
aggregating two steps gives, relative to its reference,
$N:1.067\pm0.018$, $K_{\rm tot}:1.158\pm0.030$ and
$Y:1.118\pm0.024$. Aggregating four steps, the production ratio
is $1.346\pm0.026$. The protocol starts from an empty state,
uses three seeds and compares $1200<t_{\rm old}\le2000$;
flows are divided by the aggregation factor, not the stocks.
The per-seed values and other variants are in
\chemin{data_article/temps_graines.csv} and
\chemin{data_article/temps_points.csv}. These trials and the counterexample
establish the absence of an exact symmetry; they do not constitute
a study of convergence toward a continuous limit.

